\section{Síntesis y ampliación}

\subsection{Hilo conductor central de esta clase}

Esta clase parte del perceptrón más básico y, siguiendo una ruta de evolución clara, construye gradualmente una comprensión completa del modelado de secuencias:

\begin{enumerate}
    \item \textbf{Perceptrón/red feedforward}: puede procesar la entrada de un solo paso de tiempo, pero no puede aprovechar la información histórica
    \item \textbf{RNN}: introduce una ``memoria'' mediante el estado oculto, con lo que logra la transferencia de información entre pasos de tiempo
    \item \textbf{LSTM}: mitiga el problema del desvanecimiento del gradiente mediante un mecanismo de compuertas, mejorando el aprendizaje de dependencias de largo plazo
    \item \textbf{Attention/Transformer}: abandona por completo la estructura recurrente y modela directamente las dependencias globales mediante el mecanismo de autoatención
\end{enumerate}

\begin{importantbox}{El cambio de paradigma de RNN a Transformer}
Esta ruta de evolución refleja un cambio importante de filosofía de diseño en el aprendizaje profundo:
\begin{itemize}
    \item \textbf{Filosofía de RNN}: procesar en orden temporal, manteniendo un estado de memoria comprimido
    \item \textbf{Filosofía de Transformer}: ver todo el contexto de una vez y dejar que la red decida por sí misma en qué enfocarse
\end{itemize}
El éxito de Transformer muestra que, en lugar de que los humanos diseñen la forma en que fluye la información (como la recurrencia), es mejor dejar que la red descubra por sí misma, mediante el aprendizaje, los patrones importantes de los datos.
\end{importantbox}

\subsection{Referencia rápida de conceptos clave}

\begin{table}[H]
\centering
\caption{Tabla de referencia rápida de conceptos clave de esta clase}
\begin{tabular}{ll}
\toprule
\textbf{Concepto} & \textbf{Punto clave} \\
\midrule
Hidden State ($h_t$) & Memoria interna que la RNN transmite entre pasos de tiempo \\
Recurrence Relation & $h_t = f_W(x_t, h_{t-1})$, actualiza de forma recurrente el estado oculto \\
BPTT & Algoritmo que retropropaga el gradiente a lo largo del eje temporal \\
Vanishing Gradient & El gradiente se atenúa con los pasos de tiempo, lo que dificulta el aprendizaje de dependencias de largo plazo \\
LSTM & Variante de RNN que protege el flujo del gradiente mediante un mecanismo de compuertas \\
Embedding & Representación aprendible que mapea símbolos discretos a vectores continuos \\
Positional Encoding & Inyecta información de posición en la entrada del Transformer \\
Self-Attention & Cálculo de atención Query-Key-Value dentro de la secuencia \\
Scaled Dot-Product & $\text{softmax}(QK^T/\sqrt{d_k})V$ \\
Multi-Head Attention & Múltiples grupos de Attention en paralelo, que capturan diversos patrones de atención \\
Transformer & Arquitectura de modelo de secuencias basada por completo en Attention \\
\bottomrule
\end{tabular}
\end{table}

\subsection{Lecturas adicionales}

\begin{itemize}
    \item Vaswani et al., \textit{Attention Is All You Need}, NeurIPS 2017. Artículo pionero de Transformer.
    \item Hochreiter \& Schmidhuber, \textit{Long Short-Term Memory}, Neural Computation 1997. Artículo original de LSTM.
    \item Devlin et al., \textit{BERT: Pre-training of Deep Bidirectional Transformers for Language Understanding}, NAACL 2019.
    \item Brown et al., \textit{Language Models are Few-Shot Learners (GPT-3)}, NeurIPS 2020.
    \item Dosovitskiy et al., \textit{An Image is Worth 16x16 Words: Transformers for Image Recognition at Scale (ViT)}, ICLR 2021.
    \item Sitio oficial del curso MIT 6.S191: \url{http://introtodeeplearning.com}
    \item MIT 6.S191 Lab 1: Deep Learning in Python and Music Generation with RNNs
\end{itemize}

%% --- Fin del contenido --- %%

\end{document}
